\scanpage{4}
\mspage{2}
que le $\mathbb Q$-type d'homotopie de $X$ est récupéré si $X$ est simplement connexe -- de façon plus précise, il y a \inserted{(sauf erreur)} une équivalence de cat. \inserted{donnée par $C^{\ast}_{\mathrm{DRS}}$} entre la cat. hom. \inserted{faible} des espaces connexes et simplement connexes (du pt de vue singulier), et une certaine catégorie dérivée formée avec les $\mathbb Q$-algèbres graduées associatives et \inserted{anti}comm. \deleted{telles} à degrés $\geq 0$ telles que
\[
H^0(A)\simeq\mathbb Q\ \deleted{\unsure{H^1(A)\simeq 0}}\quad H^1(A)\simeq 0 .
\]
Il y a une théorie des modèles minimaux \inserted{pour de telles algèbres}, une façon très simple de récupérer les $\pi_i\otimes_{\mathbb Z}\mathbb Q$ en termes d'un tel modèle.--- (On renvoie au papier de Sullivan)

Mais à nouveau, on perd la torsion !

\noindent\rule{\linewidth}{0.3pt}

Dém. du th. d'isomorphie de Sullivan
\[
H^\bullet(C^\bullet_{\mathrm{DRS}\deleted{\unsure{\mathbb Q}}}(X,\mathbb Q))\simeq H^\bullet(X,\mathbb Q)
\]

$\forall n\geq 0$, $[n]$ simplexe type de dim $n$
\[
\left(\sum_0^n X_i=1,\ \deleted{\unsure{\varepsilon}}X_i\geq 0\right)
\]
$\Delta^{[n]}\deleted{\unsure{\mathbb R}}$ sa réalisation géométrique

$E^{[n]}$ l'espace affine qui le sous tend (avec une $\mathbb Q$-structure)
\margin{$\displaystyle\left|\begin{array}{l}\sum_0^n X_i=1\\ \text{dans }\mathbb R^{n+1}\end{array}\right.$}

\[
\mathrm{DRS}^{\bullet}_{\deleted{\unsure{i}}[n]}=C^\bullet_{\mathbb Q\text{-}\mathrm{DR}}(E^{[n]})\quad\text{son complexe de De Rham $\mathbb Q$-algébrique}
\]
C'est contravariant en $[n]$, d'où
\[
\mathrm{DRS}^{\bullet\deleted{\unsure{\ast}}}_{\ast}=(\mathrm{DRS}_{[n]})_{n\geq 0}
\]
Algèbre différentielle graduée \deleted{(semi)}\inserted{semi-}simpliciale (et $=$ simplicial) -- à degrés $\geq 0$, anticommutative.

Pour tt espace topologique $X$, $S_{\ast}(X)$ son complexe singulier $1/2$ simplicial. On a
\[
C^\bullet_{\mathrm{DRS}}(X,\mathbb Q)\simeq\mathrm{Hom}(S_{\unsure{\ast}}(X),\mathrm{DRS}^{\bullet}_{\ast})
\]

% archive sheet number: 3
